% TOP_PROBLEM: {"id": 10000038, "title": "Cheeger constant and the percolation nonuniqueness phase", "category_id": 19, "category": {"id": 19, "name": "probability", "display_name": "Probability", "description": "Problems involving probability theory, stochastic processes, and random structures.", "slug": "probability", "order_index": 19, "created_at": "2026-07-31T15:26:25.671Z"}, "status": "partially_solved", "problem_number": "AMR-099-0038", "set_id": 13}
% FIRST_POSTED: 2026-10-01
% ENTRY_KIND: statement_only
% UnsolvedMath ID 10000038; source catalogue code AMR-099-0038
% !TeX program = lualatex
% Definition and sources only; this file is not an LLM solution attempt.
\documentclass[11pt,a4paper]{article}
\usepackage[margin=25mm]{geometry}
\usepackage{fontspec}
\setmainfont{lmroman10-regular.otf}[BoldFont=lmroman10-bold.otf,
 ItalicFont=lmroman10-italic.otf,BoldItalicFont=lmroman10-bolditalic.otf]
\usepackage{amsmath,amssymb,mathtools}
\usepackage{unicode-math}
\setmathfont{latinmodern-math.otf}
\AtBeginDocument{\let\setminus\smallsetminus}
\usepackage{xurl,hyperref}
\hypersetup{colorlinks=true,urlcolor=blue}
\setcounter{secnumdepth}{0}
\setlength{\parindent}{0pt}
\setlength{\parskip}{5pt}
\setlength{\emergencystretch}{3em}
\begin{document}
\section*{Cheeger constant and the percolation nonuniqueness phase}\label{10000038}
{\small UnsolvedMath ID 10000038 · AMR-099-0038 · Probability · AMR Open Problem Lists}
\subsection{Definitions and mathematical statement}
Let $G$ be an infinite connected locally finite vertex-transitive graph. For each $p\in[0,1]$, perform Bernoulli bond percolation by retaining edges independently with probability $p$, and let $N_\infty(p)$ be the number of infinite open components. Define
\[
p_c(G)=\inf\{p:\mathbb P_p(N_\infty(p)\ge1)>0\},
\qquad
p_u(G)=\inf\{p:\mathbb P_p(N_\infty(p)=1)=1\}.
\]
For finite nonempty $A\subset V(G)$, write $\partial_G A$ for the external vertex boundary, and set $h(G)=\inf_A|\partial_G A|/|A|$.

Prove the equivalence
\[
p_c(G)<p_u(G)\quad\Longleftrightarrow\quad h(G)>0.
\]
In particular, the difficult direction asks whether every nonamenable transitive graph has a nonempty interval of parameters on which there are infinitely many infinite open clusters almost surely. The constants and the length of this interval may depend on the graph.

Local finiteness and transitivity imply a common finite degree. No restriction to Cayley graphs, unimodular graphs, planar graphs, or hyperbolic graphs is part of the assertion. The definition of $p_u$ specifies uniqueness with probability one, rather than uniqueness conditional on a selected root being in an infinite cluster. The target makes no separate assertion about what happens at either endpoint $p_c$ or $p_u$. It compares the existence of an intermediate percolation phase with expansion measured by all finite vertex sets in the underlying graph.
\subsection{Short English statement}
Does every nonamenable transitive graph have a percolation phase with infinitely many infinite clusters, and only such graphs?
\subsection{Sources}
\begin{itemize}
\item[S1] ULAM.AI and the UnsolvedMath Contributors, supplied problems.json, record 10000038 (AMR-099-0038), AMR Open Problem Lists. Catalogue identity and source transcription; CC BY 4.0.
\url{https://huggingface.co/datasets/ulamai/UnsolvedMath}.
\item[S2] Benjamini - Coarse Geometry and Randomness (Saint-Flour notes), Conjecture 9.10, PDF page 65. Original source indexed by AMR-099-0038.
\url{https://www.wisdom.weizmann.ac.il/~itai/stflouraug24.pdf#page=65}.
\item[S3] I. Benjamini, Coarse Geometry and Randomness, primary Saint-Flour notes; the numbered question and its surrounding definitions.
\url{https://arquivo.pt/noFrame/replay/20201231041548id_/http://www.wisdom.weizmann.ac.il/~itai/stflouraug24.pdf}.
\end{itemize}
\end{document}
